Goiás e o Distrito Federal (\textasciitilde340.000~km\textsuperscript{2} do bioma Cerrado)
concentram a tensão central do planejamento do uso da terra nos trópicos: expandir e intensificar
a produção sem comprometer os ecossistemas e os recursos hídricos dos quais ela depende. Esta
dissertação constrói, para GO~e~DF a 250~m, um \textbf{atlas aberto de viabilidade agrícola
multissegmento com zoneamento e projeção climática}, abrangendo sete segmentos de uso da terra:
soja, cana-de-açúcar, outras culturas anuais, piscicultura, pecuária, conservação de
floresta/Cerrado e geração fotovoltaica. A viabilidade atual resulta de uma avaliação multicritério
nos moldes da FAO (pertinência fuzzy e pesos AHP). O zoneamento agroambiental vem de um agrupamento
k-means não supervisionado sobre uma representação descorrelacionada das variáveis. E a projeção
para meados do século aplica o método dos fatores de mudança a um conjunto de cinco modelos
CMIP6. Todo o processamento é executado no lado do servidor, no Google Earth Engine, \textbf{sem
upload de dados externos}, com validação contra o uso atual da terra (MapBiomas) e um indicador de
produtividade (MODIS MOD17). Os resultados respondem às três questões de pesquisa: superfícies de
viabilidade atual para os sete segmentos e uma partição em \textbf{dez zonas} com vocações
comparativas distintas, entre elas um planalto alto cuja melhor vocação é a conservação. As zonas,
construídas apenas com variáveis biofísicas, reproduzem a geografia observada do uso da terra sem
usar a cobertura como insumo (AUC da soja de 0,871). Elas também revelam um hiato relevante entre
viabilidade potencial e uso atual, sobretudo em pastagens, apontando para intensificação --- não
expansão sobre vegetação nativa --- como direção mais provável de mudança do uso da terra na
região. A mudança climática de meados do século
é modesta,
porém direcionalmente robusta (concordância entre modelos $\approx$ 1,0 no teste mais rigoroso):
atinge sobretudo a
conservação e, em menor grau, o grupo milho/algodão/sorgo e a geração fotovoltaica, enquanto os
planaltos nobres de soja e cana permanecem essencialmente estáveis. O método roda inteiramente na
nuvem a partir de dados públicos, sem upload, o que o torna reprodutível a custo marginal baixo
para outros estados do Cerrado.

\vspace{1em}
\noindent\textbf{Palavras-chave:} viabilidade agrícola de terras; avaliação de terras (FAO);
processo analítico hierárquico (AHP); zoneamento agroambiental; Cerrado; Google Earth Engine;
projeção climática CMIP6; análise multicritério de decisão.
